\section{Peripheral Channel Structure and Transfer-Operator Gaps: Supporting Results}
\label{ch:spectral_supporting_results}

This section collects the generic supporting results for finite-dimensional
quantum channels behind the peripheral spectral theory of
Section~\ref{sec:peripheral_refinements}: an auxiliary trace identity and
trace-positivity result, the Banach-algebra convergence fact behind the
overlap-decay theorems, the rank-one Perron projection behind the primitive
overlap limit, and the cyclic-decomposition machinery that removes peripheral
periodicity.

\section{Auxiliary trace identities}
\label{sec:app_spectral_mixed_trace}


\begin{lemma}[Trace expansion over matrix units]\label{lem:trace_expansion}
    \lean{Matrix.linearMap_trace_eq_sum_apply_single}
    \leanok
    For any linear endomorphism $T$ on $M_{D_1\times D_2}(\C)$, the operator
    trace expands as
    \begin{align}
        \Tr(T)
         & =\sum_{p=0}^{D_1-1}\sum_{q=0}^{D_2-1}
                              (T(E_{pq}))_{pq},
        \label{eq:spectral_trace_expansion}
    \end{align}
    where $E_{pq}\in M_{D_1\times D_2}(\C)$ has entry $1$ in position
    $(p,q)$ and zero elsewhere.
\end{lemma}

\begin{proof}\leanok
    The operator trace is
    $\Tr(T)=\sum_{p,q}\tr(T(E_{pq})E_{qp})$, and
    $\tr(ME_{qp})=M_{pq}$.
\end{proof}


\section{Auxiliary trace positivity}

\begin{lemma}[Trace product of positive semidefinite matrices]
    \label{lem:psd_trace_product_nonneg}
    \lean{Matrix.PosSemidef.trace_mul_nonneg}
    \leanok
    If $A,B\geq0$, then $\tr(AB)\geq0$.
\end{lemma}

\begin{proof}\leanok
    Write $B=U\Lambda U^\dagger$ with $U$ unitary and $\Lambda$ diagonal.
    Then
    \begin{align}
        \tr(AB)
         & =\tr\!\left((U^\dagger AU)\Lambda\right)
        =\sum_i(U^\dagger AU)_{ii}\Lambda_{ii}
        \geq0,
        \notag
    \end{align}
    because $U^\dagger AU\geq0$ and $\Lambda_{ii}\geq0$ for every $i$.
\end{proof}

\section{Spectral-radius decay and overlap limits}
\label{sec:app_spectral_decay}

\begin{lemma}[Powers vanish below unit spectral radius]%
    \label{lem:spectral_radius_pow_zero}
    \lean{pow_tendsto_zero_of_spectralRadius_lt_one}
    \leanok
    Let $a$ be an element of a complex Banach algebra with
    $\rho_{\spec}(a)<1$. Then $a^n\to0$ as $n\to\infty$.
\end{lemma}

\begin{proof}
    \leanok
    Choose $r$ with $\rho_{\spec}(a)<r<1$. The Gelfand formula
    $\|a^n\|^{1/n}\to\rho_{\spec}(a)$ gives $\|a^n\|\leq r^n$ for all
    sufficiently large $n$, and $r^n\to0$.
\end{proof}

\section{Rank-one Perron projection}
\label{sec:app_spectral_perron_projection}

\begin{align}
    P(X)
     & =\frac{\tr(X)}{\tr(\rho)}\rho,
    \notag
\end{align}
and $N:=E-P$ for the complementary part. The power decomposition
(Theorem~\ref{thm:pow_decomp}, Chapter~\ref{ch:channels}) gives
$E^n=P+N^n$ for $n\geq1$.

\begin{lemma}[Trace of the fixed-point projection]\label{lem:fixedPointProj_trace}
    \lean{fixedPointProj_trace}
    \leanok
    \uses{def:fixed_point_proj}
    The fixed-point projection $P$ has operator trace $\Tr(P)=1$.
\end{lemma}

\begin{proof}\leanok
    $P$ has the rank-one form $P(X)=f(X)\rho$ for the linear functional
    $f(X)=\tr(X)/\tr(\rho)$. The operator trace of a rank-one endomorphism
    $X\mapsto f(X)\rho$ equals $f(\rho)$, and
    $f(\rho)=\tr(\rho)/\tr(\rho)=1$.
\end{proof}

\begin{theorem}[Complementary gap at a prescribed fixed point]%
    \label{thm:primitive_fixedPoint_compl_lt_one_irred_channel}
    \lean{spectralRadius_compl_lt_one_of_primitive_fixedPoint_of_irreducible_channel}
    \leanok
    \uses{def:fixed_point_proj, def:irreducible_cp, def:primitive_channel}
    Let $E:\MN{D}\to\MN{D}$ be an irreducible primitive channel, and let
    $\rho\ne0$ be a positive semidefinite fixed point. Then
    $\tr(\rho)\ne0$ and, for $P=P_\rho$,
    $\rho_{\spec}(E-P)<1$.
\end{theorem}

\begin{proof}\leanok
    \uses{thm:primitive_compl_lt_one,thm:spectralRadius_lt_one_of_eigenvalues_lt_one}
    Positive semidefiniteness and $\rho\ne0$ imply $\tr(\rho)\ne0$.
    Every eigenvalue of $E-P$ has modulus strictly less than one by
    irreducibility and primitivity. Since $\MN{D}$ is finite-dimensional,
    the maximum eigenvalue modulus equals the spectral radius.
\end{proof}

\section{Periodicity removal}
\label{sec:app_spectral_periodicity}

This section proves the cyclic-decomposition machinery of
Theorem~\ref{thm:cyclic_decomposition_irreducible_schwarz} and the
divisibility statements of Theorem~\ref{thm:peripheral_cyclic_group} and
Theorem~\ref{thm:channel_period_divides_dim}: normalizing a peripheral
eigenvector to a unitary, forming its discrete Fourier (cyclic) spectral
projections, and tracking how the channel shifts the resulting corners.

\begin{lemma}[Peripheral eigenvectors are invertible]%
    \label{lem:isUnit_peripheral_eigenvector}
    \lean{Kraus.isUnit_peripheral_eigenvector}
    \leanok
    \uses{def:kraus_map, def:unital_kraus, def:kraus_adjoint_map,
        def:irreducible_cp}
    Let $E$ be an irreducible unital Kraus map on $\MN{D}$ with a
    positive-definite adjoint fixed point.
    If $X \ne 0$ satisfies $E(X) = \mu X$ with $|\mu| = 1$,
    then $X$ is invertible.
\end{lemma}

\begin{proof}
    \leanok
    \uses{thm:ks_equality_of_peripheral_eigenvector_of_fixedPoint,
        thm:psd_fp_pd_irred}
    By the Kadison--Schwarz inequality
    (\ref{eq:schwarz_kadison}), the gap
    $G=E(X^\dagger X)-E(X)^\dagger E(X)$ is positive semidefinite.
    Pairing $G$ with the positive-definite adjoint fixed point and using
    $E(X)=\mu X$ and $|\mu|=1$ gives weighted trace zero, hence $G=0$.
    Therefore $E(X^\dagger X)=E(X)^\dagger E(X)=X^\dagger X$, so
    $X^\dagger X$ is a nonzero positive semidefinite fixed point.
    Irreducibility upgrades it to a positive-definite, hence invertible,
    matrix, so $\det(X)\ne 0$.
\end{proof}

\begin{lemma}[Peripheral eigenvalues: product closure]%
    \label{lem:peripheral_mul_closure}
    \lean{Kraus.peripheralEigenvalues_mul_mem_of_irreducible_unital_of_adjoint_fixedPoint}
    \leanok
    \uses{def:kraus_map, def:unital_kraus, def:kraus_adjoint_map,
        def:irreducible_cp, def:peripheral_eigenvalues}
    Let $E$ be an irreducible unital Kraus map on $\MN{D}$ with a
    positive-definite adjoint fixed point.
    If $\mu, \nu$ are peripheral eigenvalues of $E$, then so is $\mu\nu$.
\end{lemma}

\begin{proof}
    \leanok
    \uses{lem:isUnit_peripheral_eigenvector,
        thm:ks_equality_of_peripheral_eigenvector_of_fixedPoint,
        thm:right_multiplicative_identity}
    Take eigenvectors $X$ and $Y$ for $\mu$ and $\nu$, respectively.
    The Kadison--Schwarz equality at $X$ places $X$ in the multiplicative
    domain. Thus the right multiplicative identity
    (\ref{eq:schwarz_right_multiplicative}) gives
    $E(YX)=E(Y)E(X)=(\nu\mu)YX$.
    Since $X$ and $Y$ are invertible by
    Lemma~\ref{lem:isUnit_peripheral_eigenvector}, $YX\ne 0$, and
    $|\mu\nu|=1$.
\end{proof}

\begin{lemma}[Scalar fixed points for irreducible unital maps]
    \label{lem:irreducible_unital_fixed_points_scalar}
    \lean{Kraus.fixed_eq_scalar_of_irreducible_unital}
    \leanok
    \uses{thm:psd_fp_unique_irred_cp}
    If $E$ is irreducible and unital, then every fixed point of $E$ is a
    scalar multiple of $\Id$.
\end{lemma}

\begin{proof}
    \leanok
    \uses{thm:psd_fp_unique_irred_cp}
    If $E(X)=X$, then the Hermitian matrices
    $X+X^\dagger$ and $\mathrm{i}(X-X^\dagger)$ are also fixed by $E$.
    The positive-semidefinite fixed-point uniqueness theorem gives
    $X+X^\dagger=a\Id$ and
    $\mathrm{i}(X-X^\dagger)=b\Id$, hence
    $X=\frac{1}{2}(a-\mathrm{i}b)\Id$.
\end{proof}

\begin{lemma}[Peripheral unitary eigenvector]
    \label{lem:peripheral_unitary_eigenvector}
    \lean{Kraus.exists_peripheral_unitary_of_irreducible_schwarz}
    \leanok
    \uses{thm:ks_equality_of_peripheral_eigenvector_of_fixedPoint,
        thm:psd_fp_unique_irred_cp}
    For an irreducible unital Schwarz map with faithful adjoint fixed point,
    each peripheral eigenvalue admits a unitary eigenvector.
\end{lemma}

\begin{proof}\leanok
    For an eigenvector $X$ at a peripheral eigenvalue $\gamma$, the
    Kadison--Schwarz equality gives
    $E(X^\dagger X)=E(X)^\dagger E(X)=X^\dagger X$, so $X^\dagger X$ is a
    nonzero positive semidefinite fixed point. Irreducible uniqueness of the
    positive semidefinite fixed point makes it a positive scalar $c\Id$;
    rescaling $X$ by $c^{-1/2}$ produces a unitary eigenvector.
\end{proof}

\begin{lemma}[Powers on a peripheral-unitary orbit]
    \label{lem:peripheral_unitary_powers}
    \lean{Kraus.map_powers_of_peripheral_unitary}
    \leanok
    \uses{thm:ks_equality_of_peripheral_eigenvector_of_fixedPoint,
        thm:right_multiplicative_identity}
    If $E(U)=\mu U$ with $|\mu|=1$ and $U$ unitary, then
    $E(U^n)=\mu^nU^n$ for every $n$.
\end{lemma}

\begin{proof}\leanok
    The Kadison--Schwarz equality at the unitary eigenvector places $U$ in
    the multiplicative domain, so the right multiplicative identity
    (\ref{eq:schwarz_right_multiplicative}) gives
    $E(U^n)=E(U)^n=\mu^nU^n$ by induction on $n$.
\end{proof}

\begin{lemma}[Normalized peripheral unitary]
    \label{lem:normalized_peripheral_unitary}
    \lean{Kraus.exists_normalized_peripheral_unitary_of_irreducible_schwarz}
    \leanok
    \uses{lem:peripheral_unitary_eigenvector, lem:peripheral_unitary_powers,
        lem:irreducible_unital_fixed_points_scalar}
    Peripheral unitary eigenvectors can be normalized compatibly with
    the faithful fixed-point state.
\end{lemma}

\begin{proof}\leanok
    Given a unitary eigenvector $U$ for a primitive $m$-th root $\gamma$,
    Lemma~\ref{lem:peripheral_unitary_powers} gives $E(U^m)=U^m$, a fixed
    point that is scalar by
    Lemma~\ref{lem:irreducible_unital_fixed_points_scalar}: $U^m=\alpha\Id$.
    Unitarity of $U^m$ forces $|\alpha|=1$; multiplying $U$ by an $m$-th
    root $\beta$ of $\alpha^{-1}$ produces $(\beta U)^m=1$ while preserving
    the eigenvector equation.
\end{proof}

\begin{lemma}[Cyclic projections from a peripheral unitary]
    \label{lem:cyclic_projections_from_peripheral_unitary}
    \lean{exists_cyclic_projections_of_peripheral_unitary}
    \leanok
    A primitive $m$-th peripheral root yields $m$ orthogonal projections
    summing to $\Id$ and cyclically permuted by the channel.
\end{lemma}

\begin{proof}\leanok
    \uses{lem:normalized_peripheral_unitary}
    Given a unitary $U$ of order $m$ with $E(U)=\gamma U$ for a primitive
    $m$-th root $\gamma$, the Fourier projections
    $P_k:=m^{-1}\sum_{j=0}^{m-1}\overline\gamma^{kj}U^j$ are mutually
    orthogonal and sum to $\Id$ by discrete Fourier orthogonality of the
    characters $j\mapsto\gamma^{kj}$; since $E(U^j)=\gamma^jU^j$
    (Lemma~\ref{lem:peripheral_unitary_powers}), each Fourier mode is
    shifted by one step, giving $E(P_{k+1})=P_k$.
\end{proof}

Combining these lemmas proves Theorem~\ref{thm:cyclic_decomposition_irreducible_schwarz}:
Lemma~\ref{lem:normalized_peripheral_unitary} supplies a unitary generator of
exact order $m$ for a primitive $m$-th peripheral root, and
Lemma~\ref{lem:cyclic_projections_from_peripheral_unitary} builds the $m$
Fourier projections out of its powers, using
Lemma~\ref{lem:peripheral_unitary_powers} to identify how the channel shifts
them.

\subsection{Cyclic corners and restricted primitivity}

\begin{definition}[Corner preservation]
    \lean{PreservesCorner}
    \leanok
    For an orthogonal projection $P$ and linear map $E$, the corner
    $P\MN{D}P$ is preserved if $E(PXP)=PE(PXP)P$ for every $X$.
\end{definition}

\begin{definition}[Corner subspace]
    \lean{cornerSubmodule}
    \leanok
    The corner subspace associated with $P$ is
    $\{PXP:X\in\MN{D}\}$.
\end{definition}

\begin{definition}[Restricted corner map]
    \label{def:corner_restriction}
    \lean{cornerRestriction}
    \leanok
    The restriction of $E$ to an invariant corner $P\MN{D}P$.
\end{definition}

\begin{definition}[Irreducible restriction on a corner]
    \lean{IsIrreducibleOnCorner}
    \leanok
    Irreducibility of the corner-restricted map.
\end{definition}

\begin{definition}[Rank of an orthogonal projection]
    \lean{cornerRank}
    \leanok
    The rank of an orthogonal projection $P\in\MN{D}$, equivalently the
    dimension of its range.
\end{definition}

\begin{lemma}[Compression isometry for an orthogonal projection]%
    \label{thm:compression_isometry_projection}
    \lean{cornerSubmoduleMatrixLinearEquiv}
    \leanok
    Let $P\in\MN{D}$ be an orthogonal projection of rank $n$. The corner
    algebra $P\MN{D}P$ is linearly isomorphic to the full matrix algebra
    $\MN{n}$. The isomorphism is built from the spectral diagonalisation
    of $P$ by conjugating the top-left $n\times n$ block by the unitary
    diagonalising $P$.
\end{lemma}

\begin{lemma}[Rank equals trace for an orthogonal projection]%
    \label{thm:corner_rank_eq_trace}
    \lean{cornerRank_eq_trace}
    \leanok
    The rank of an orthogonal projection $P$ equals its trace:
    $n=\tr(P)$.
\end{lemma}

\begin{proof}\leanok
    The eigenvalues of an orthogonal (Hermitian idempotent) projection are
    $0$ or $1$: the eigenvalues $\lambda$ satisfy $\lambda^2=\lambda$ from
    $P^2=P$. The trace is the sum of the eigenvalues, which counts the
    eigenvalue-$1$ multiplicities, i.e.\ the rank $n$.
\end{proof}

\begin{lemma}[Power maps preserve each cyclic sector]
    \label{lem:corner_pow_preserved}
    \lean{preserves_corner_pow_of_cyclic_decomp}
    \leanok
    Appropriate powers of the channel preserve each sector corner.
\end{lemma}

\begin{proof}\leanok
    Induction on the exponent, using the left- and right-multiplicative
    domain identities on each sector projection $P_k$, shows
    $T^n(P_{k+n\bmod m}XP_{k+n\bmod m})=P_k\,T^n(X)\,P_k$; specializing to
    $n=m$, where $P_{k+m}=P_k$, gives that $T^m$ preserves the corner
    $P_k\MN{D}P_k$.
\end{proof}

\begin{lemma}[Primitivity passes to a fixed corner]
    \label{lem:primitivity_corner_restriction}
    \lean{isPrimitive_cornerRestriction_of_isPrimitive}
    \leanok
    \uses{def:corner_restriction}
    Let $P\in\MN{D}$ be a nonzero idempotent and let $E$ be a linear map on
    $\MN{D}$ that preserves the corner $P\MN{D}P$, is primitive, and fixes
    the corner projection, $E(P)=P$. Then the restriction of $E$ to the
    corner $P\MN{D}P$ is again primitive.
\end{lemma}

\begin{proof}
    \leanok
    \uses{def:corner_restriction}
    The corner projection $P$ lies in its own corner and is fixed by the
    restriction, so $1$ is a peripheral eigenvalue of the restriction.
    Conversely, every eigenvalue of the restriction of unit modulus arises
    from a nonzero corner element $X$ with $E(X)=\mu X$; viewing $X$ as an
    element of $\MN{D}$ exhibits $\mu$ as a peripheral eigenvalue of $E$.
    Primitivity of $E$ forces $\mu=1$, so the peripheral spectrum of the
    restriction is exactly $\{1\}$ and the restriction is primitive.
\end{proof}

\begin{lemma}[Irreducibility of restricted sector dynamics]
    \label{lem:corner_restriction_irreducible}
    \lean{isIrreducible_restriction_of_cyclic_decomp}
    \leanok
    Let $T$ be irreducible, let $P_0,\ldots,P_{m-1}$ be orthogonal
    projections summing to $\Id$ and cyclically permuted by $T$, and
    suppose that every orthogonal projection $Q$ with $QP_k=P_kQ=Q$ that
    is invariant under the corner restriction of $T^m$ to $P_k\MN{D}P_k$
    admits an orthogonal projection $R$ invariant under $T$ on the full
    algebra with $Q=0\iff R=0$ and $Q=P_k\iff R=\Id$. Then, for every $k$,
    the restriction of $T^m$ to the corner $P_k\MN{D}P_k$ is irreducible.
\end{lemma}

\begin{proof}\leanok
    Let $Q$ be an invariant projection for the corner restriction of $T^m$
    on $P_k\MN{D}P_k$. The orbit-sum hypothesis supplies an ambient
    invariant projection $R$ for $T$ with $Q=0\iff R=0$ and
    $Q=P_k\iff R=\Id$. Irreducibility of $T$ forces $R=0$ or $R=\Id$, hence
    $Q=0$ or $Q=P_k$ by the two biconditionals.
\end{proof}

\begin{lemma}[Primitivity of restricted sector dynamics]
    \label{lem:corner_restriction_primitive}
    \lean{isPrimitive_restriction_of_cyclic_decomp}
    \leanok
    Each sector restriction is primitive.
\end{lemma}

\begin{proof}\leanok
    \uses{lem:primitivity_corner_restriction}
    Every peripheral eigenvalue of $T$ is an $m$-th root of unity, so the
    peripheral spectrum of $T^m$ is exactly $\{1\}$; combined with
    $T^m(P_k)=P_k$, Lemma~\ref{lem:primitivity_corner_restriction} makes the
    corner restriction of $T^m$ to $P_k\MN{D}P_k$ primitive.
\end{proof}

\begin{lemma}[Corner invariance at the period]
    \label{lem:corner_invariance_perm}
    \lean{preserves_corner_pow_orderOf_of_perm_decomp}
    \leanok
    The channel raised to the period preserves every cyclic corner.
\end{lemma}

\begin{proof}\leanok
    The same induction as in Lemma~\ref{lem:corner_pow_preserved}, with the
    cyclic shift $k\mapsto k+1$ replaced by a permutation $\sigma$ of the
    sector index set, shows that $T^{\mathrm{ord}(\sigma)}$ preserves each
    corner $P_k\MN{D}P_k$.
\end{proof}

\subsection{Group structure and divisibility}

\begin{lemma}[Peripheral product closure]
    \label{lem:peripheral_closed_under_mul_group}
    \lean{PeripheralSpectrum.peripheral_eigenvalues_closed_under_mul}
    \leanok
    \uses{def:kraus_map, def:unital_kraus, def:kraus_adjoint_map,
        def:irreducible_cp, def:peripheral_eigenvalues}
    The product of two peripheral eigenvalues is peripheral.
\end{lemma}

\begin{proof}\leanok
    \uses{lem:peripheral_unitary_eigenvector,
        thm:ks_equality_of_peripheral_eigenvector_of_fixedPoint,
        thm:right_multiplicative_identity}
    Peripheral eigenvalues admit unitary eigenvectors
    (Lemma~\ref{lem:peripheral_unitary_eigenvector}). If $U_\alpha,U_\beta$
    are unitary eigenvectors for $\alpha,\beta$, the multiplicative-domain
    identity gives
    $E(U_\alpha U_\beta)=E(U_\alpha)E(U_\beta)=\alpha\beta\,U_\alpha U_\beta$,
    and $U_\alpha U_\beta\neq0$ since both factors are unitary.
\end{proof}

\begin{lemma}[Peripheral inverse closure]
    \label{lem:peripheral_closed_under_inv_group}
    \lean{PeripheralSpectrum.peripheral_eigenvalues_closed_under_inv}
    \leanok
    \uses{def:kraus_map, def:unital_kraus, def:kraus_adjoint_map,
        def:irreducible_cp, def:peripheral_eigenvalues}
    The inverse of a peripheral eigenvalue is peripheral.
\end{lemma}

\begin{proof}\leanok
    \uses{lem:peripheral_unitary_eigenvector}
    For a unitary eigenvector $U_\alpha$ with $|\alpha|=1$, conjugating the
    eigenvector equation gives
    $E(U_\alpha^\dagger)=\overline\alpha\,U_\alpha^\dagger
        =\alpha^{-1}U_\alpha^\dagger$, and $U_\alpha^\dagger\neq0$.
\end{proof}

Theorem~\ref{thm:peripheral_multiplicity_one} follows the same pattern:
given a peripheral unitary eigenvector $U$ for $\gamma$
(Lemma~\ref{lem:peripheral_unitary_eigenvector}), the multiplicative-domain
identity gives $E(XU^\dagger)=XU^\dagger$ for any $\gamma$-eigenvector $X$,
so the scalar fixed-point lemma for irreducible unital maps
(Lemma~\ref{lem:irreducible_unital_fixed_points_scalar}) gives
$XU^\dagger=c\Id$, hence $X=cU$ and the eigenspace is one-dimensional.

The divisibility statements of Theorem~\ref{thm:peripheral_cyclic_group} and
Theorem~\ref{thm:channel_period_divides_dim} follow from
Lemma~\ref{lem:cyclic_projections_from_peripheral_unitary} and
Theorem~\ref{thm:corner_rank_eq_trace}: the $m$ cyclic Fourier projections
$P_0,\ldots,P_{m-1}$ sum to $\Id$ and are cyclically permuted by $E$, so
trace preservation gives $\tr(P_{k+1})=\tr(P_k)$ for every $k$, hence
$D=\tr(\Id)=\sum_{k=0}^{m-1}\tr(P_k)=m\,\tr(P_0)$. Since the trace of an
orthogonal projection is a natural number (Theorem~\ref{thm:corner_rank_eq_trace}),
$m\mid D$.
